\section{Représenter un raisonnement}
\label{sec:representation}

\subsection{Visualiser des arguments}

\textcite{ref09} formulent le problème central : saisir la structure hiérarchique d’un argument tout en pouvant lire le texte de ses prémisses pose des problèmes d’aperçu, de détail et de navigation. À partir d’entretiens avec des analystes d’arguments, ils ont conçu trois techniques qui combinent une visualisation d’arbre (\emph{sunburst} ou \emph{icicle}) et un affichage du contenu, intégré à la figure ou placé dans une infobulle. Les analystes jugeaient le \emph{sunburst} bon pour l’aperçu et préféraient le texte intégré ; une étude contrôlée avec 21 participants et trois tâches désigne pourtant le \emph{sunburst} avec infobulle comme le meilleur compromis. La tension entre aperçu et lecture n’a donc pas de solution unique : elle se règle selon la tâche. Ces arguments sont des arbres, ce qui autorise des représentations par remplissage de l’espace ; les nôtres sont des graphes avec des liens transversaux et des contradictions, d’où le choix d’une représentation nœuds-liens.

\subsection{Du texte à la carte avec des LLM}

Trois systèmes récents associent extraction par LLM et visualisation interactive. Story Ribbons \parencite{ref01} trace les trajectoires des personnages et des thèmes à plusieurs niveaux narratifs pour 36 œuvres. Ses trois motifs d’interaction, les explications à la demande, les dimensions définies en langage naturel et les questions posées au modèle, sont conçus pour pallier les limites des systèmes fondés sur l’IA \parencite[1, 5]{ref01}. Parmi les 16 participants de l’étude, 14 disent pourtant que les comportements inattendus du modèle n’ont pas entamé leur confiance dans les visualisations \parencite[8]{ref01}. Une erreur visible ne rend donc pas le lecteur prudent de lui-même : ce qui n’a pas été vérifié doit être signalé comme tel.

Graphologue \parencite{ref02} transforme les réponses de GPT-4 en diagrammes nœuds-liens construits en direct ; le survol d’un nœud surligne le passage correspondant de la réponse, et inversement \parencite[8]{ref02}. Son étude utilisateur porte sur sept personnes \parencite[2]{ref02}. Sensecape \parencite{ref03} organise l’exploration d’un sujet en niveaux d’abstraction, avec un zoom sémantique et une vue hiérarchique. Dans une étude intra-sujets à 12 participants, ceux-ci explorent davantage de concepts qu’avec une interface de référence associant conversation et toile simple (68,3 contre 22,8 en moyenne) \parencite[9--10]{ref03}. Les auteurs ont volontairement omis l’indication des sources, pour ne pas distraire les participants, et écrivent qu’un déploiement réel exigerait des mécanismes pour vérifier ce que produit le modèle \parencite[14]{ref03}.

La revue de \textcite{ref04} élargit ce constat. Sur 48 systèmes qui associent LLM et visualisation, 25 ont fait l’objet d’une évaluation avec des utilisateurs \parencite[17]{ref04}, et peu d’articles ont mesuré ou cherché à améliorer la fiabilité des informations produites par le modèle. Les auteurs y voient un obstacle à l’adoption de ces outils et plaident pour un cadre d’évaluation qui prenne en compte la confiance \parencite[20]{ref04}.

\subsection{Principes d’interaction}

Le principe formulé par \textcite[336]{ref27}, d’abord l’aperçu, puis le zoom et le filtre, enfin les détails à la demande, décrit presque terme à terme l’organisation du prototype. L’aperçu montre cinq à douze étapes principales, et chaque étape se déplie en sous-étapes. Des cases filtrent les types de relations, et une option estompe tout ce qui n’appartient pas à la chaîne de l’étape choisie. Le panneau de détails présente les phrases citées, l’extrait dans sa mise en page d’origine, les termes et les remarques. La taxonomie de Shneiderman comporte aussi la tâche \enquote{relier}, qui correspond ici au passage d’une étape à ce qui l’appuie.

\subsection{Disposer le graphe}

Pour placer les étapes, le prototype utilise ELK Layered, qui met en œuvre l’approche en couches de \textcite{ref10} : les nœuds sont affectés à des couches, les couches sont réordonnées pour réduire les croisements, puis les coordonnées sont calculées \parencite{ref31}. Cette approche oriente le plus d’arêtes possible dans une même direction, ce qui convient à une lecture qui va des raisons vers les conclusions. Les contradictions n’entrent pas dans cet ordre. L’ordre du texte n’est pas perdu pour autant : la lecture guidée peut suivre l’ordre des phrases, et la page du texte intégral situe chaque étape principale dans une marge.

Le sujet évoque aussi des représentations de type \emph{storyline}, qui montrent l’évolution d’entités au fil d’un récit. StoryFlow formule leur disposition comme une optimisation hybride, discrète puis continue, qui réduit les croisements, les ondulations et l’espace perdu assez vite pour permettre l’interaction en temps réel \parencite[2436]{ref11a}. iStoryline montre que ces critères géométriques ne suffisent pas à produire ce que des personnes dessineraient à la main, et intègre des interactions de l’utilisateur dans l’optimisation \parencite[769]{ref11b}. Un axe logique n’est cependant pas une chronologie : le prototype garde la disposition en couches pour la carte et réserve l’ordre du texte à la lecture guidée et à la marge du texte intégral. Une vue de type \emph{storyline}, où chaque étape serait placée selon sa position dans le texte, reste à construire.

Le regroupement d’arêtes (\emph{edge bundling}) simplifie le dessin des graphes denses en rapprochant les arêtes semblables \parencite{ref12b}. La méthode de \textcite{ref12} construit une carte de densité des arêtes par estimation à noyau, puis déplace les arêtes le long du gradient de cette densité. Avec cinq à douze étapes, chaque relation reste lisible sans regroupement, et chacune doit pouvoir être vérifiée isolément. Le regroupement deviendrait utile si l’on dépliait de nombreuses sous-étapes à la fois ou si l’on superposait les cartes de plusieurs documents.
